\documentclass[12pt]{article}
\usepackage[utf8]{inputenc}
\usepackage[spanish]{babel}
\usepackage{amsmath}
\usepackage{amsfonts}
\usepackage{amssymb}
\usepackage{geometry}
\geometry{letterpaper, margin=2.5cm}

\begin{document}

\begin{center}
    \Large \textbf{Evaluación de Examen 1: Unidad I} \\
    \large Física para Ciencias de la Salud (005-1713) \\
    \vspace{0.5cm}
    \textbf{Estudiante:} Yilselys Naranjo Betancourt \quad \textbf{C.I.:} 33.385.584
    \vspace{0.3cm} \\
    \large \textbf{Calificación Final: 9/10 puntos}
\end{center}

\vspace{0.5cm}

\section*{Análisis de Resultados}

\subsection*{1. Ángulo plano (Conversión a radianes)}
\textbf{Procedimiento y resultado:} Plantea correctamente la conversión. A pesar de omitir un cero al inicio $\left(\frac{\pi}{18}\right)$, lo corrige en la siguiente igualdad $\left(\frac{75\pi}{180^\circ}\right)$. Simplifica la fracción paso a paso hasta obtener la expresión exacta e irreductible de $\frac{5\pi}{12}$ rad. \\
\textbf{Veredicto:} \textbf{Correcto} (2/2 pts).

\subsection*{2. Sistemas de coordenadas (Longitud del hueso)}
\textbf{Procedimiento y resultado:} Utiliza la ecuación formal de la distancia entre dos puntos y sustituye las coordenadas de manera precisa. Desarrolla aritméticamente las restas, los binomios al cuadrado y la suma ($\sqrt{36+64} = \sqrt{100}$) para concluir con el resultado exacto de $10$ cm. \\
\textbf{Veredicto:} \textbf{Correcto} (2/2 pts).

\subsection*{3. Vectores (Fuerza resultante)}
\textbf{Procedimiento y resultado:} La estudiante efectúa las sumas algebraicas respetando la ley de los signos para obtener 30 en el eje $\hat{i}$ y 60 en el eje $\hat{j}$. Existe un leve detalle de notación al denotar estas sumas parciales como $\vec{F}_1$ y $\vec{F}_2$ (nombres que ya pertenecían a los vectores originales), pero ensambla la resultante final de forma impecable: $\vec{F}_R = (30\hat{i} + 60\hat{j})$ N. \\
\textbf{Veredicto:} \textbf{Correcto} (2/2 pts).

\subsection*{4. Potencias de diez (Notación científica y prefijo)}
\textbf{Procedimiento y resultado:} Transforma la cifra decimal de forma matemáticamente válida empleando potencias de base diez ($12,5 \times 10^{-6}$ metros). Sin embargo, omite la respuesta final específica requerida en el enunciado: identificar y nombrar el prefijo del Sistema Internacional (micrómetros, $\mu$m). \\
\textbf{Veredicto:} \textbf{Incompleto / Parcialmente correcto} (1/2 pts).

\subsection*{5. Sistemas de unidades (Conversión de densidad)}
\textbf{Procedimiento y resultado:} Aplica los factores de conversión en cadena de forma correcta. Expresa las relaciones fraccionarias exactas para masa y volumen, simplificando la ecuación para llegar a la operación directa ($1,03 \cdot 1000$). Obtiene el resultado numérico y de unidades totalmente exacto: $1030 \text{ Kg/m}^3$. \\
\textbf{Veredicto:} \textbf{Correcto} (2/2 pts).

\vspace{0.5cm}
\section*{Comentarios Generales y Sugerencias}
¡Excelente examen! La estudiante muestra un nivel analítico muy sólido, manteniendo un buen orden procedimental a la hora de estructurar sus operaciones matemáticas y la simplificación de fracciones.
\begin{itemize}
    \item Se recomienda cuidar el rigor de la notación en álgebra vectorial (nombrar a las componentes rectangulares como $\vec{F}_{Rx}$ y $\vec{F}_{Ry}$ en lugar de reutilizar los nombres $\vec{F}_1$ y $\vec{F}_2$ en la Pregunta 3).
\end{itemize}

\end{document}